\selectlanguage{ngerman}%
\frenchspacing%
\begin{ledgroupsized}[r]{120mm}%
\footnotesize%
\pstart%
\noindent%
\textbf{Überlieferung:}%
\pend%
\end{ledgroupsized}%
\begin{ledgroupsized}[r]{114mm}%
\footnotesize%
\pstart%
\parindent -6mm%
\makebox[6mm][l]{\textit{E}}%
Druck nach unbekannter Vorlage:
\textsc{Gerland} 1906, S.\;11 (N.\;8, titellos).%
\cite{00197}
Bei der verschollenen handschriftlichen Vorlage soll es sich laut \textit{E} um ein \textit{Kleines schlecht geschriebenes Blättchen von Leibnizens Hand} gehandelt haben.
Der oft fehlerhafte Text von \textit{E} (Gerlands Transkription) wird hier wiedergegeben und gegebenenfalls berichtigt oder kommentiert; die Rechtschreibung wird an die Editionsregeln angepasst.
\pend%
\end{ledgroupsized}%
\vspace{5mm}%
\begin{ledgroup}%
\footnotesize%
\pstart%
\noindent%
\textbf{Datierungsgründe:}
Im titellosen und undatierten Text N.\;\ref{56218}, der nach heutigem Wissensstand lediglich im posthumen Druck \textit{E} %
überliefert ist, wird auf die Beschreibung des Karpaten-Gebirges\indexO{Karpaten} in Daniel Fröhlichs\indexN{\textls{Frölich} (Frolichius, Frölichius), David 1595\textendash1648} %
\textit{Bibliotheca seu Cynosura peregrinantium}\cite{00047} %
hingewiesen (siehe unten den Erläuterungsapparat).
Ob Leibniz diese Beschreibung aus erster Hand kannte oder aber aus deren Wiedergabe in Otto von Guerickes\indexN{\textls{Guericke} (Gerickius, Gericke), Otto von 1602\textendash1686} %
\textit{Experimenta nova} (siehe unten den Erläuterungsapparat), %
es steht fest, dass Erwähnungen von Frölichs Werk im Leibniz-Nachlss nach heutigem Wissensstand frühestens auf das Jahr 1672 zurückgehen (vgl. \textit{LSB} VIII,\;1 N.\;36, S.\;263.6--7, 266.3; \textit{LSB} VIII,\;1 N.\;62, S.\;570, Randbemerkung zu Z.\;15\:f.).\cite{01197}\cite{01278}
Es ist demgemäß unwahrscheinlich, dass die Aufzeichung N.\;\ref{56218} vor 1672 verfasst wurde.
\pend%
\pstart%
An Fröhlichs\indexN{\textls{Frölich} (Frolichius, Frölichius), David 1595\textendash1648} Beschreibung knüpft Leibniz erneut in dem auf die Zeitspanne 1681 bis 1682 datierbaren Entwurf \textit{Explicatio mechanica soni} an (\textit{LSB} VIII,\;3 N. 12\textsubscript{2}, S.\;110.18).\cite{01481}
Da keine weiteren Erwähnungen dieses Berichts im Leibniz-Nachlass bekannt sind, dürfte auch der Text N.\;\ref{56218} nicht viel später entstanden sein.
Die Aufzeichnung N.\;\ref{56218} könnte tatsächlich noch Anfang der 1680er Jahre entstanden sein, als Leibniz sich intensiv mit Fragen der physikalischen Akustik, insbesondere mit der Entstehung und Übertragung des Schalls, befasste.
\pend%
\pstart%
Bei seinem Bericht über ein besonderes Widerhallphänomen in der Helmstedter\indexO{Helmstedt} Umgebung, das für die Studenten der dortigen \textit{Academia Julia} eine Vergnügungsangelegenheit gewesen sein soll, bezieht sich Leibniz im Text N.\;\ref{56218} auf eine (nicht ermittelte) Quelle, die offenbar auch Otto von Guericke\indexN{\textls{Guericke} (Gerickius, Gericke), Otto von 1602\textendash1686} bekannt war (siehe unten den Erläuterungsapparat).
Auch dieser Umstand verfestigt den Eindruck, dass der Text N.\;\ref{56218} eher in frühen Jahren entstand, d.\,h. nicht viele Jahre später als 1672.
\pend%
\end{ledgroup}%
\selectlanguage{latin}%
\frenchspacing%
\newpage%
\normalsize%
\pstart%
\noindent%
\lbrack S.\:11\rbrack\
Der Donner%
\indexS{Donner}
scheint großenteils ein Echo%
\indexS{Echo}
in den Wolken%
\indexS{Wolken}
zu seyn,
wie denn solcher aus der
\edtext{%
relation des
\edtext{%
F\lbrack r\rbrack olichii%
\indexN{\textls{Frölich} (Frolichius, Frölichius), David 1595\textendash1648}%
}{%
\lemma{Fiolichii}%
\Bfootnote{%
\textit{L\;ändert Hrsg.}%
}}
de Carpathis monte%
\indexO{Karpaten}%
}{%
\lemma{relation \lbrack...\rbrack\ monte}%
\Cfootnote{%
Siehe die Beschreibung der bei der Besteigung des Karpaten-Gebirges\indexO{Karpaten} beobachteten Phänomene des Widerhalls in
\textsc{D.\;Frölich}, %
\textit{Bibliotheca seu Cynosura peregrinantium}, %
\glqq Geographia\grqq, cap.\;IX, §\;372\,f.
(Band II.1, Ulm 1643, S.\;287\,f.).\cite{00047} %
Vom natürlichen Donner ist dort allerdings nirgendwo die Rede.
Zitiert war Frölichs%
\indexN{\textls{Frölich} (Frolichius, Frölichius), David 1595\textendash1648}
Beschreibung auch in
\textsc{O.\;v.\;Guericke}, \textit{Experimenta nova},
lib.\;V, cap.\;8 (Amsterdam 1672, S.\;162).\cite{00055}
Das Zitat wurde von Leibniz in seinen Guericke-Auszügen exzerpiert
(\textit{LSB} VIII,\;1 N.\;36, S.\;266.5–7).\cite{01197}%
}}
und anderen umbständen abzunehmen.
Daher auch der Donner%
\indexS{Donner}
fortzulauffen und sich uns zu nähern scheint.
\edtext{%
Etwa eine Meile%
\indexS{Meile}
von
\edtext{%
Helmstad%
\indexO{Helmstedt}%
}{%
\lemma{Helmstad}%
\Cfootnote{%
Die Hansestadt Helmstedt war von 1576 bis 1810 Sitz der protestantischen Universität \textit{Academia Julia}.%
}}
bey einem Dorff%
\indexS{Dorf}
\edtext{%
zum Drewdel%
\indexO{Drewdel (?) bei Helmstedt}%
}{%
\lemma{zum Drewdel}%
\Cfootnote{%
Nicht ermittelt.%
}}
ist ein
\edtext{%
aus%
}{%
\lemma{aus}%
\Cfootnote{%
Die Lesung ist unsicher.%
}}
der maßen schöhnes Echo,%
\indexS{Echo}
welches den Donner%
\indexS{Donner}
wohl nachahmet.
Die Studenten%
\indexS{Student}
kommen bisweilen dahin,
bringen kleine mörser%
\indexS{Mörser}
mit hinauß
und schießen im Holz%
\indexS{Holz}
herumb%
\lbrack;\rbrack\
in der ebene%
\indexS{Ebene}
zeiget sich wie ein
\edtext{%
gelbes(?)%
}{%
\lemma{gelbes(?)}%
\Cfootnote{%
Die Lesung ist, wie \textit{E} selbst vermerkt, unsicher.%
}}
theatrum,%
\indexS{theatrum}
der schall%
\indexS{Schall}
antwortet erstlich sehr stark ein baar mahl,
dann wird er immer geschwinder und schwächer
und laufft also wie ein Donner%
\indexS{Donner}
an dem Wald%
\indexS{Wald}
herumb.%
}{%
\lemma{Etwa \bracks{...} herumb}%
\Cfootnote{%
Das Phänomen eines besonderen Widerhalls in der Helmstedter Umgebung\indexO{Helmstedt} findet Erwähnung auch in
\textsc{Guericke}, \textit{Experimenta nova},
lib.\;IV, cap.\;10 (S.\;139b).\cite{00055}
Dies kann allerdings nicht die (einzige) Quelle von Leibnizens detaillirterem Bericht sein.%
}}
\pend%
\selectlanguage{latin}%
\frenchspacing%
\normalsize%
\newpage%
